\section{Conclusion}
\label{sec:conclusion}

We presented a cross-metric depth$\times$width sweep of cross-layer representation similarity for pure ReLU MLPs on \cifar and \mnist. MLP hidden states at different layers are neither identical nor unrelated: mean off-diagonal \ckam sits firmly inside $(0,1)$, with adjacent layers consistently more similar than far-apart layers. Depth drives the emergence of block structure---at $d\!=\!16$, a contiguous central block of layers collapses to CKA $\gtrsim 0.9$, reproducing in pure MLPs a phenomenon previously seen only in CNNs, ResNets, and transformers. Width, at low-to-moderate depth, has the \emph{opposite} effect to its role in cross-\emph{model} CNN comparisons: wider MLPs differentiate their layers more, directly confirming that limited per-layer transformation capacity is what forces narrow MLPs to re-use features across depth.

\para{Key takeaway.} The motivating intuition is correct and quantifiable: MLP hidden states are similar enough that every internal layer carries most of the information of its neighbors, yet different enough that depth, width, and regime leave distinct, measurable signatures in the similarity structure.

\para{Future work.} We flag four directions. (i)~\textbf{Training trajectories}: record \ckam every epoch to see \emph{when} block structure forms, extending the bottom-up convergence analysis of \citet{raghu2017svcca} to MLPs. (ii)~\textbf{Residual MLPs}: adding skip connections should flatten cross-layer CKA~\citep{raghu2021vit}; directly test this in MLPs. (iii)~\textbf{Pruning}: prune the block-interior layers of a $d\!=\!16$ MLP and measure the accuracy cost. Our interpretation predicts near-zero loss. (iv)~\textbf{Normalization}: \citet{daneshmand2021bn} shows BatchNorm orthogonalizes hidden features; the effect on observed block structure remains open for MLPs.
